% ═══════════════════════════════════════════════════════════════
% MUSTERLÖSUNG: Waagerechter Wurf (Physik Klasse 10E)
% Version 3.0 – Passend zu AB-01 Version 3.0
% ═══════════════════════════════════════════════════════════════

\documentclass[10pt, a4paper]{article}

\usepackage[utf8]{inputenc}
\usepackage[T1]{fontenc}
\usepackage[ngerman]{babel}
\usepackage{helvet}
\renewcommand{\familydefault}{\sfdefault}

\usepackage[margin=1.5cm, top=1.8cm, bottom=1.2cm]{geometry}
\usepackage{enumitem}
\usepackage{tabularx}
\usepackage{array}
\usepackage{multicol}
\usepackage{amsmath}
\usepackage{amssymb}

\usepackage{xcolor}
\usepackage{tcolorbox}
\usepackage{fancyhdr}
\usepackage{lastpage}
\usepackage{tikz}

\setlist{nosep, leftmargin=1.5em}

% ═══════════════════════════════════════════════════════════════
% FARBEN
% ═══════════════════════════════════════════════════════════════

\definecolor{physik}{HTML}{2c5aa0}
\definecolor{grau}{HTML}{888888}
\definecolor{hellgrau}{HTML}{f5f5f5}
\definecolor{loesung}{HTML}{c62828}

\newcommand{\fachfarbe}{physik}

% ═══════════════════════════════════════════════════════════════
% VARIABLEN
% ═══════════════════════════════════════════════════════════════

\newcommand{\thema}{Waagerechter Wurf}
\newcommand{\fach}{Physik}
\newcommand{\klasse}{10E}
\newcommand{\gesamt}{26}

% ═══════════════════════════════════════════════════════════════
% KOPF- UND FUSSZEILE
% ═══════════════════════════════════════════════════════════════

\pagestyle{fancy}
\fancyhf{}
\renewcommand{\headrulewidth}{0.4pt}
\renewcommand{\footrulewidth}{0pt}

\lhead{\small\fach{} \klasse{} -- \thema{} \textcolor{loesung}{\textbf{[MUSTERLÖSUNG]}}}
\rhead{\small Name: \rule{4cm}{0.4pt}}

\cfoot{\small\thepage/\pageref{LastPage}}

% ═══════════════════════════════════════════════════════════════
% BEFEHLE
% ═══════════════════════════════════════════════════════════════

\newcommand{\luecke}[1]{\rule{#1}{0.4pt}}
\newcommand{\ph}[1]{\textcolor{\fachfarbe}{\textbf{#1}}}
\newcommand{\punkte}[1]{\hfill\small\textcolor{grau}{/#1}}
\newcommand{\lsg}[1]{\textcolor{loesung}{\textbf{#1}}}

\newtcolorbox{merkbox}[1][]{
    colback=hellgrau,
    colframe=\fachfarbe,
    boxrule=0.8pt,
    arc=0pt,
    left=5pt, right=5pt, top=5pt, bottom=5pt,
    fonttitle=\bfseries\small,
    title=#1
}

\newtcolorbox{formelbox}{
    colback=white,
    colframe=\fachfarbe,
    boxrule=1.2pt,
    arc=0pt,
    left=6pt, right=6pt, top=6pt, bottom=6pt
}

\newcommand{\abschnitt}[1]{%
    \vspace{4pt}
    \noindent\textbf{\textcolor{\fachfarbe}{#1}}
    \vspace{2pt}
}

% ═══════════════════════════════════════════════════════════════
% DOKUMENT
% ═══════════════════════════════════════════════════════════════

\begin{document}

\begin{center}
{\large\bfseries Arbeitsblatt: \thema}\\[2pt]
{\small\textcolor{loesung}{MUSTERLÖSUNG}}
\end{center}

\vspace{-4pt}
\hrule
\vspace{4pt}

% ═══════════════════════════════════════════════════════════════
% SEITE 1
% ═══════════════════════════════════════════════════════════════

\begin{multicols}{2}
\small

% ─────────────────────────────────────────────────────────────
% K1: Überlagerungsprinzip
% ─────────────────────────────────────────────────────────────

\abschnitt{K1: Das Überlagerungsprinzip}

\begin{merkbox}[Merksatz]
Beim waagerechten Wurf überlagern sich zwei \ph{unabhängige} Bewegungen:

\vspace{4pt}
\begin{tabular}{@{}p{2cm}l@{}}
\textbf{Horizontal:} & \lsg{gleichförmige Bewegung} \\[8pt]
\textbf{Vertikal:} & \lsg{freier Fall} \\
\end{tabular}

\vspace{6pt}
Die horizontale Geschwindigkeit beeinflusst die \lsg{Fallzeit} \textbf{nicht}.

\vspace{4pt}
Deshalb landen Fallball und Wurfball \lsg{gleichzeitig}.
\end{merkbox}

\vspace{2pt}

% ─────────────────────────────────────────────────────────────
% K2: Formeln
% ─────────────────────────────────────────────────────────────

\abschnitt{K2: Formeln (aus LP übertragen)}

\begin{formelbox}
\textbf{Fallzeit:}
\vspace{4pt}

$t = $ \lsg{$\sqrt{\dfrac{2h}{g}}$}

\vspace{2pt}
\footnotesize
$t$ = Fallzeit [s], $h$ = Höhe [m], $g = 10\,\text{m/s}^2$
\end{formelbox}

\vspace{3pt}

\begin{formelbox}
\textbf{Wurfweite:}
\vspace{4pt}

$s_x = $ \lsg{$v_0 \cdot t$}

\vspace{2pt}
\footnotesize
$s_x$ = Wurfweite [m], $v_0$ = Anfangsgeschw. [m/s]
\end{formelbox}

\vspace{2pt}

% ─────────────────────────────────────────────────────────────
% K3: Aufprall
% ─────────────────────────────────────────────────────────────

\abschnitt{K3: Aufprallgeschwindigkeit \& -winkel}

\begin{formelbox}
\textbf{Vertikale Endgeschwindigkeit:}
\vspace{4pt}

$v_y = $ \lsg{$g \cdot t$}

\vspace{6pt}
\textbf{Aufprallgeschwindigkeit} (Pythagoras):
\vspace{4pt}

$v_{\text{res}} = $ \lsg{$\sqrt{v_x^2 + v_y^2}$}

\vspace{6pt}
\textbf{Aufprallwinkel} (Trigonometrie):
\vspace{4pt}

$\tan(\alpha) = $ \lsg{$\dfrac{v_y}{v_x}$}
\end{formelbox}

\vspace{2pt}

% ─────────────────────────────────────────────────────────────
% Tan-Tabelle
% ─────────────────────────────────────────────────────────────

\abschnitt{Tangens-Tabelle}

\begin{center}
\begin{tabular}{|c|c|c|c|c|c|}
\hline
$\alpha$ & 30° & 37° & 45° & 53° & 60° \\
\hline
$\tan(\alpha)$ & 0,58 & 0,75 & 1,00 & 1,33 & 1,73 \\
\hline
\end{tabular}
\end{center}

% ─────────────────────────────────────────────────────────────
% SPALTENUMBRUCH
% ─────────────────────────────────────────────────────────────

\columnbreak

% ─────────────────────────────────────────────────────────────
% Ü1: Grundaufgabe
% ─────────────────────────────────────────────────────────────

\abschnitt{Ü1: Grundaufgabe}\punkte{4}

Ein Ball wird aus $h = 20\,\text{m}$ Höhe mit $v_0 = 15\,\text{m/s}$ waagerecht geworfen.

\textbf{a)} \textbf{Berechne} die Fallzeit $t$ des Balls.\hfill\textcolor{grau}{\small/2}

\lsg{$t = \sqrt{\frac{2 \cdot 20}{10}} = \sqrt{4} = 2\,\text{s}$}

\vspace{0.4cm}

\textbf{b)} \textbf{Berechne} die Wurfweite $s_x$.\hfill\textcolor{grau}{\small/2}

\lsg{$s_x = 15 \cdot 2 = 30\,\text{m}$}

\vspace{0.4cm}

\hrule
\vspace{4pt}

% ─────────────────────────────────────────────────────────────
% Ü2: Skizze Vektoren
% ─────────────────────────────────────────────────────────────

\abschnitt{Ü2: Geschwindigkeitsvektoren}\punkte{2}

\textbf{Zeichne} am mittleren Ball die Vektoren $\vec{v}_x$, $\vec{v}_y$ und $\vec{v}_{\text{res}}$ maßstabsgerecht ein.

\begin{center}
\begin{tikzpicture}[scale=0.52]
    % Koordinatensystem
    \draw[->] (0,0) -- (7,0) node[right] {\small$x$};
    \draw[->] (0,0) -- (0,5) node[above] {\small$y$};

    % Boden
    \draw[thick] (-0.3,0) -- (7,0);

    % Plattform
    \fill[gray!40] (-0.3,4.2) rectangle (0.8,4.5);
    \draw[thick] (-0.3,4.2) -- (0.8,4.2);

    % Höhe h
    \draw[<->, thick] (-0.5,0) -- (-0.5,4.2);
    \node[left] at (-0.5,2.1) {\small$h$};

    % Parabel (gestrichelt)
    \draw[dashed, gray, thick] (0.5,4.2) parabola (6,0);

    % Ball Start
    \fill[\fachfarbe] (0.5,4.2) circle (0.12);

    % Ball Mitte (AUF der Parabel: x=3 → y≈3.33)
    \fill[\fachfarbe!50] (3,3.33) circle (0.12);

    % Ball Ende
    \fill[\fachfarbe] (6,0) circle (0.12);

    % Wurfweite
    \draw[<->, thick] (0.5,-0.6) -- (6,-0.6);
    \node[below] at (3.25,-0.6) {\small$s_x$};

    % v0 Pfeil
    \draw[->, very thick, red!70!black] (0.5,4.2) -- (2.2,4.2) node[above] {\small$v_0$};

    % LÖSUNGS-VEKTOREN am mittleren Ball
    \draw[->, very thick, loesung] (3,3.33) -- (4.5,3.33) node[above] {\small\lsg{$\vec{v}_x$}};
    \draw[->, very thick, loesung] (3,3.33) -- (3,2.03) node[left] {\small\lsg{$\vec{v}_y$}};
    \draw[->, very thick, loesung] (3,3.33) -- (4.2,2.23) node[right] {\small\lsg{$\vec{v}_{\text{res}}$}};
\end{tikzpicture}
\end{center}

\hrule
\vspace{4pt}

% ─────────────────────────────────────────────────────────────
% Ü3: Aufprall (war Ü2)
% ─────────────────────────────────────────────────────────────

\abschnitt{Ü3: Aufprallwerte}\punkte{6}

Verwende die Werte aus Ü1 ($h = 20\,\text{m}$, $v_0 = 15\,\text{m/s}$).

\textbf{a)} \textbf{Berechne} die vertikale Geschwindigkeit $v_y$ beim Aufprall.\hfill\textcolor{grau}{\small/2}

\lsg{$v_y = 10 \cdot 2 = 20\,\text{m/s}$}

\vspace{0.4cm}

\textbf{b)} \textbf{Berechne} die resultierende Aufprallgeschwindigkeit $v_{\text{res}}$.\hfill\textcolor{grau}{\small/2}

\lsg{$v_{\text{res}} = \sqrt{15^2 + 20^2} = \sqrt{625} = 25\,\text{m/s}$}

\vspace{0.4cm}

\textbf{c)} \textbf{Bestimme} den Aufprallwinkel $\alpha$ mithilfe der Tangens-Tabelle.\hfill\textcolor{grau}{\small/2}

\lsg{$\tan(\alpha) = \frac{20}{15} = 1{,}33 \Rightarrow \alpha \approx 53°$}

\vspace{0.2cm}

\hrule
\vspace{4pt}

% ─────────────────────────────────────────────────────────────
% Ü4: Anwendung (war Ü3)
% ─────────────────────────────────────────────────────────────

\abschnitt{Ü4: Anwendung -- Hilfspaket}\punkte{10}

Ein Flugzeug fliegt in $h = 80\,\text{m}$ Höhe mit $v_0 = 30\,\text{m/s}$ und wirft ein Hilfspaket ab. \textbf{Berechne} die folgenden Größen:

\textbf{a)} Fallzeit $t = $ \lsg{$\sqrt{16} = 4\,\text{s}$} \hfill\textcolor{grau}{\small/2}

\textbf{b)} Wurfweite $s_x = $ \lsg{$30 \cdot 4 = 120\,\text{m}$} \hfill\textcolor{grau}{\small/2}

\textbf{c)} Vertikale Geschwindigkeit $v_y = $ \lsg{$10 \cdot 4 = 40\,\text{m/s}$} \hfill\textcolor{grau}{\small/2}

\textbf{d)} Aufprallgeschwindigkeit $v_{\text{res}} = $ \lsg{$\sqrt{900 + 1600} = 50\,\text{m/s}$} \hfill\textcolor{grau}{\small/2}

\textbf{e)} Aufprallwinkel $\alpha = $ \lsg{$\tan^{-1}(1{,}33) = 53°$} \hfill\textcolor{grau}{\small/2}

\vspace{0.1cm}
\hrule
\vspace{4pt}

% ─────────────────────────────────────────────────────────────
% Ü5: Umkehraufgabe (war Ü4)
% ─────────────────────────────────────────────────────────────

\abschnitt{Ü5: Umkehraufgabe}\punkte{4}

Ein Paket fällt $t = 4\,\text{s}$ und landet $s_x = 100\,\text{m}$ entfernt.

\textbf{a)} \textbf{Berechne} die Anfangsgeschwindigkeit $v_0$ des Flugzeugs.\hfill\textcolor{grau}{\small/2}

\lsg{$v_0 = \frac{s_x}{t} = \frac{100}{4} = 25\,\text{m/s}$}

\vspace{0.3cm}

\textbf{b)} \textbf{Berechne} die Abwurfhöhe $h$.\hfill\textcolor{grau}{\small/2}

\lsg{$h = \frac{1}{2} g t^2 = \frac{1}{2} \cdot 10 \cdot 16 = 80\,\text{m}$}

\vspace{0.2cm}

\end{multicols}

% ═══════════════════════════════════════════════════════════════
% SEITE 2
% ═══════════════════════════════════════════════════════════════

\newpage

\begin{multicols}{2}
\small

% ─────────────────────────────────────────────────────────────
% K4: Schräger Wurf
% ─────────────────────────────────────────────────────────────

\abschnitt{K4: Schräger Wurf (qualitativ)}

\begin{merkbox}
Der schräge Wurf ist eine Überlagerung von:

\vspace{4pt}
\begin{tabular}{@{}rl@{}}
1. & gleichförmige Horizontalbewegung \\[6pt]
2. & senkrechter Wurf nach oben \\
\end{tabular}

\vspace{8pt}
Bei symmetrischem Wurf gilt: Steigzeit = Fallzeit

\vspace{8pt}
Maximale Wurfweite bei Winkel: \lsg{45°}

\vspace{8pt}
Bahnform: Wurfparabel
\end{merkbox}

\vspace{4pt}

% ─────────────────────────────────────────────────────────────
% K5: Orbitalmechanik
% ─────────────────────────────────────────────────────────────

\abschnitt{K5: Orbitalmechanik}

\begin{merkbox}
Ein Orbit ist ein kontinuierlicher \textbf{waagerechter Wurf}.

Der Satellit „fällt" ständig, aber die Erde krümmt sich genauso schnell weg.

\vspace{8pt}
\textbf{Erste kosmische Geschwindigkeit:}

$v_1 \approx$ \lsg{7,9} km/s \quad (Kreisbahn um die Erde)

\vspace{6pt}
\textbf{Zweite kosmische Geschwindigkeit:}

$v_2 \approx$ \lsg{11,2} km/s \quad (Flucht aus Erdanziehung)
\end{merkbox}

\columnbreak

% ─────────────────────────────────────────────────────────────
% Fallzeiten-Tabelle
% ─────────────────────────────────────────────────────────────

\abschnitt{Typische Fallzeiten ($g = 10\,\text{m/s}^2$)}

\begin{center}
\begin{tabular}{|c|c|c|}
\hline
\textbf{Höhe $h$} & \textbf{$\frac{2h}{g}$} & \textbf{Fallzeit $t$} \\
\hline
5 m & 1 & 1 s \\
\hline
20 m & 4 & 2 s \\
\hline
45 m & 9 & 3 s \\
\hline
80 m & 16 & 4 s \\
\hline
125 m & 25 & 5 s \\
\hline
\end{tabular}
\end{center}

\vspace{4pt}

% ─────────────────────────────────────────────────────────────
% Pythagoräische Tripel
% ─────────────────────────────────────────────────────────────

\abschnitt{Nützlich: Pythagoräische Tripel}

\begin{center}
\begin{tabular}{|c|c|c|}
\hline
$a$ & $b$ & $c = \sqrt{a^2 + b^2}$ \\
\hline
3 & 4 & 5 \\
\hline
5 & 12 & 13 \\
\hline
15 & 20 & 25 \\
\hline
30 & 40 & 50 \\
\hline
\end{tabular}
\end{center}

\end{multicols}

% ═══════════════════════════════════════════════════════════════
% SELBSTCHECK
% ═══════════════════════════════════════════════════════════════

\vspace{-2pt}
\hrule
\vspace{3pt}

\small
\textbf{Selbstcheck:} \hfill
$\boxtimes$ Überlagerungsprinzip \quad
$\boxtimes$ Fallzeit ($t$) \quad
$\boxtimes$ Wurfweite ($s_x$) \quad
$\boxtimes$ Aufprallgeschw. ($v_{\text{res}}$) \quad
$\boxtimes$ Aufprallwinkel ($\alpha$)

\vspace{3pt}
\hrule
\vspace{3pt}

\begin{tabular}{llll}
\textbf{Gesamt:} & \lsg{26} / \gesamt{} BE &
\textbf{NP:} & \lsg{15 NP}
\end{tabular}

\end{document}
